\paragraph{Mesopotamia (Sumerians, Akkadians, Babylonians).}
Clay tablets preserve entire “dream books” (for example, the series \textit{Iškar Zaqīqu}), lists of omens in the form “if a man dreams X, then Y.”  
Gods and heroes—from the \textit{Epic of Gilgamesh} onward—receive dreams as divine messages.  
These corpora demonstrate an institutionalized competence in dream classification, combining omenology, medicine, and political consultation \citep{oppenheim1956dreams}.

\paragraph{Egypt}
Egyptian “dream papyri” combine ritual prescription and prognostication; dreaming was used to guide both healing and political action \citep{oppenheim1956dreams}.  

\paragraph{Greece and the temple sleep of incubation}
At the healing sanctuaries of Asclepius (Epidaurus, Cos, and others), supplicants underwent ritual purification and \textit{enkoimesis}—sacred sleep within the temple.  
The stone inscriptions (\textit{iamata}) record miraculous healings and the divine instructions received in dreams \citep{lidonnici1995epidaurus,edelstein1957asclepius}.  
At the same time, the literary and technical tradition culminated in Artemidorus’ \textit{Oneirocritica} (2nd century CE), a systematic manual that merged empirical case studies with popular semantics and divinatory theory \citep{harrismccoy2012artemidorus}.
